% ============================================================
% MCM/ICM 美赛论文 LaTeX 模板
% 适配 2026 年新规：25 页上限（含摘要和附录）
% 使用 XeLaTeX 编译
% ============================================================

\documentclass[12pt]{article}

% ========== 字体和编码 ==========
\usepackage{fontspec}
\setmainfont{Times New Roman}
\usepackage[utf8]{inputenc}

% ========== 页面设置 ==========
\usepackage[letterpaper, margin=1in]{geometry}
\usepackage{fancyhdr}
\usepackage{lastpage}
\pagestyle{fancy}
\fancyhf{}
\fancyhead[L]{Team \# 2412345}   % 替换为你的队号
\fancyhead[R]{MCM/ICM 20XX}
\fancyfoot[C]{\thepage\ of \pageref{LastPage}}

% ========== 数学公式 ==========
\usepackage{amsmath, amssymb, amsthm}
\usepackage{mathtools}
\newtheorem{theorem}{Theorem}[section]
\newtheorem{lemma}[theorem]{Lemma}
\newtheorem{definition}{Definition}[section]

% ========== 图表 ==========
\usepackage{graphicx}
\usepackage{float}
\usepackage{subcaption}
\usepackage{booktabs}
\usepackage{longtable}
\usepackage{array}
\usepackage{xcolor}

% ========== 超链接和引用 ==========
\usepackage[hidelinks]{hyperref}
\usepackage{url}

% ========== 算法伪代码 ==========
\usepackage{algorithm}
\usepackage{algpseudocode}

% ========== 其他 ==========
\usepackage{enumitem}
\usepackage{multirow}
\usepackage{rotating}

\setlength{\parskip}{6pt}
\setlength{\parindent}{0pt}

% ========== 标题信息 ==========
\title{\textbf{Your Paper Title Here}}
\author{Team \# 2412345}
\date{February 20XX}

\begin{document}

% ============================================================
% 摘要页（单独占页，不计入正文，但计入 25 页限制）
% ============================================================
\maketitle
\thispagestyle{empty}

\begin{center}
\rule{\textwidth}{0.5pt}
\textbf{Summary}
\rule{\textwidth}{0.5pt}
\end{center}

\vspace{6pt}

% ---- 摘要正文 ----
% 要点：问题重述 → 建模方法 → 主要结果 → 结论
% 长度：一页以内
% 注意：摘要是最重要的部分！评审先看摘要决定是否继续读下去

Your summary goes here. A good summary includes:
\begin{enumerate}
    \item Problem restatement — what problem you solved
    \item Methodology — what models you built (1-2 sentences each)
    \item Key results — numerical results, key findings
    \item Conclusion — significance and implications
\end{enumerate}

\textbf{Keywords}: model, algorithm, sensitivity analysis, ...

\newpage

% ============================================================
% 正文
% ============================================================

\tableofcontents
\newpage

% --- 1. Introduction ---
\section{Introduction}

\subsection{Problem Background}
Background information and context of the problem.

\subsection{Problem Restatement}
Restate the problem in your own words, breaking it down into sub-problems.

\subsection{Literature Review}
Brief review of existing approaches and their limitations.

\subsection{Our Work}
Summarize what you accomplished and highlight your innovations.
Provide an overview of your modeling approach.

% --- 2. Assumptions ---
\section{Assumptions and Justifications}

\begin{itemize}[leftmargin=*]
    \item \textbf{Assumption 1}: Description. \textit{Justification}: Why this is reasonable.
    \item \textbf{Assumption 2}: Description. \textit{Justification}: Why this is reasonable.
    \item \textbf{Assumption 3}: Description. \textit{Justification}: Why this is reasonable.
\end{itemize}

% --- 3. Notations ---
\section{Notations}

\begin{table}[H]
\centering
\begin{tabular}{c|l|l}
\toprule
\textbf{Symbol} & \textbf{Definition} & \textbf{Unit} \\
\midrule
$x_i$ & Description of variable & unit \\
$y$ & Description of variable & unit \\
$N$ & Total number of samples & — \\
\bottomrule
\end{tabular}
\caption{Key notations used in this paper}
\label{tab:notations}
\end{table}

% --- 4. Model I ---
\section{Model I: [Model Name]}

\subsection{Model Establishment}

Describe the model: motivation, mathematical formulation, key equations.

\begin{equation}
    \label{eq:model1}
    f(x) = ...
\end{equation}

Explain what each term in the equation represents.

\subsection{Model Solution}

How you solved the model — analytical or numerical method.

\begin{algorithm}[H]
\caption{Algorithm for solving Model I}
\begin{algorithmic}[1]
\State \textbf{Input}: parameters $x_1, x_2, ...$
\State Step 1: Initialize parameters
\State Step 2: Compute ...
\State Step 3: Optimize ...
\State \textbf{Output}: result $y$
\end{algorithmic}
\end{algorithm}

\subsection{Results and Analysis}

Present results with figures and tables.

\begin{figure}[H]
    \centering
    % \includegraphics[width=0.7\textwidth]{figures/result1.png}
    \caption{Results of Model I}
    \label{fig:result1}
\end{figure}

% --- 5. Model II ---
\section{Model II: [Model Name]}

\subsection{Model Establishment}
...

\subsection{Model Solution}
...

\subsection{Results and Analysis}
...

% --- 6. Sensitivity Analysis (CRITICAL!) ---
\section{Sensitivity Analysis}

Sensitivity analysis is \textbf{required} for MCM/ICM papers.
It demonstrates the robustness of your model.

\subsection{Method}
Describe your sensitivity analysis method (e.g., OAT, Morris, Sobol, Latin Hypercube Sampling).

\subsection{Results}
Show how results change when key parameters vary within reasonable ranges.

\begin{figure}[H]
    \centering
    % \includegraphics[width=0.6\textwidth]{figures/sensitivity.png}
    \caption{Sensitivity analysis results: effect of parameter $\alpha$ on output}
    \label{fig:sensitivity}
\end{figure}

\subsection{Discussion}
Interpret the sensitivity results. Are there any parameters the model is particularly sensitive to? What does this imply?

% --- 7. Model Evaluation ---
\section{Model Evaluation}

\subsection{Strengths}
\begin{itemize}
    \item ...
    \item ...
\end{itemize}

\subsection{Weaknesses and Future Work}
\begin{itemize}
    \item ...
    \item ...
\end{itemize}

% --- 8. Conclusion ---
\section{Conclusion}

Summarize your work:
\begin{enumerate}
    \item What models you built
    \item What results you achieved
    \item What insights you gained
\end{enumerate}

% --- 9. References ---
\begin{thebibliography}{99}

\bibitem{ref1}
Author, A. (Year). \textit{Title of the paper}. Journal Name, Volume(Issue), Pages.

\bibitem{ref2}
Author, B. (Year). \textit{Book Title}. Publisher.

\bibitem{ref3}
COMAP. (2026). MCM/ICM 2026 Contest Rules and Instructions.
\url{https://www.comap.com/contests/mcm-icm}

\end{thebibliography}

% ============================================================
% Appendix
% ============================================================
\newpage
\appendix
\section{Code}

\subsection{Python Code — Model I}
\begin{verbatim}
# Your Python code here
import numpy as np
import pandas as pd

def model_function(x):
    return ...
\end{verbatim}

\subsection{MATLAB Code — Model II}
\begin{verbatim}
% Your MATLAB code here
function y = model_function(x)
    y = ...
end
\end{verbatim}

\section{Additional Figures and Tables}

% Additional supporting figures and tables

\section{AI Usage Report}
\textbf{AI Tools Used}: (e.g., ChatGPT, Claude, DeepSeek)

\textbf{Purpose}: (e.g., literature review assistance, code debugging, language polishing)

\textbf{Extent of Use}: Describe how AI was used and how originality was maintained.

\vspace{12pt}
\textit{All models and analyses presented in this paper are the original work of our team. AI tools were used only for [specific purposes] and all outputs were reviewed and validated by team members.}

\end{document}
